% Generated by scripts/tikz_reviews.py from the figure manifests.
\ReviewFigure{Convexity of functions and sets}
{figures/optim-smooth/convex-examples.png}
{figures/tikz/smooth-convex-examples.pdf}
{Convexity and convex sets: secants lie above a convex graph; line segments remain in a convex set.}
{Red segments distinguish failure of convexity, non-strict convexity with boundary segments, and strict convexity. The middle function has a flat interval, and the middle set has flat boundary segments.}
{smooth-convex-examples}
\ReviewFigure{Linear regression}
{figures/optim-smooth/ml-exemp-1.pdf}
{figures/tikz/smooth-linear-regression.pdf}
{Least-squares fitting of a linear predictor to training pairs.}
{The line represents the scalar-feature illustration of y=\ensuremath{\langle}a,x\ensuremath{\rangle}; point positions are schematic and are not a numerical dataset.}
{smooth-linear-regression}
\ReviewFigure{A linear classifier}
{figures/optim-smooth/ml-exemp-2.pdf}
{figures/tikz/smooth-linear-classifier.pdf}
{The parameter x is normal to the decision hyperplane in feature space.}
{Positive labels are placed on the positive-score side and negative labels on the negative-score side. The normal vector is perpendicular to the separating line.}
{smooth-linear-classifier}
\ReviewFigure{Classification losses}
{figures/optim-smooth/classif-loss.pdf}
{figures/tikz/smooth-classification-losses.pdf}
{Smooth empirical classification risk and standard margin losses.}
{The sketch labels a logistic loss divided by log(2), whereas this section defines the unscaled log(1+exp(u)); the reconstruction follows the section. The hinge is max(0,1+u). Zero margin is counted as an error. The unscaled logistic loss is not claimed to majorize the 0--1 loss.}
{smooth-classification-losses}
\ReviewFigure{Existence and uniqueness of minimizers}
{figures/optim-smooth/uniqueness.pdf}
{figures/tikz/smooth-minimizer-examples.pdf}
{An objective can have no minimizer, several minimizers, or a unique minimizer.}
{Five representative functions retain the original cases. The two no-minimizer examples distinguish an unbounded infimum from a finite unattained infimum. Vertical offsets and coordinate scales are schematic.}
{smooth-minimizer-examples}
\ReviewFigure{Coercive least squares}
{figures/optim-smooth/least-square-1.pdf}
{figures/tikz/smooth-least-squares-coercive.pdf}
{Full column rank gives a strictly convex quadratic with a unique minimizer.}
{The bowl depicts the full-column-rank case. Its height may include a nonzero residual at the minimizer; no exact fit is assumed.}
{smooth-least-squares-coercive}
\ReviewFigure{Least squares with a nontrivial kernel}
{figures/optim-smooth/least-square-2.pdf}
{figures/tikz/smooth-least-squares-flat.pdf}
{The objective is constant along translations by vectors in the kernel of A.}
{The trough is flat exactly along the kernel direction. The highlighted affine line represents x-star+ker(A), rather than an isolated minimizer.}
{smooth-least-squares-flat}
\ReviewFigure{Failure of the secant inequality}
{figures/optim-smooth/cvx-vs-noncvx-1.pdf}
{figures/tikz/smooth-secant-nonconvex.pdf}
{The graph of a convex function lies below every secant.}
{The interior point is compared with the secant at the same abscissa. The convex example also includes a supporting tangent; the nonconvex example explicitly violates the secant inequality.}
{smooth-secant-nonconvex}
\ReviewFigure{Convexity and the secant inequality}
{figures/optim-smooth/cvx-vs-noncvx-2.pdf}
{figures/tikz/smooth-secant-convex.pdf}
{The graph of a convex function lies below every secant.}
{The interior point is compared with the secant at the same abscissa. The convex example also includes a supporting tangent; the nonconvex example explicitly violates the secant inequality.}
{smooth-secant-convex}
\ReviewFigure{Strict convexity}
{figures/optim-smooth/strictly-cvx-2.pdf}
{figures/tikz/smooth-strict-convexity.pdf}
{For distinct endpoints and an interior convex combination, the secant inequality is strict.}
{The reconstruction follows the definitions and notation in the accompanying section.}
{smooth-strict-convexity}
\ReviewFigure{Convexity with an affine interval}
{figures/optim-smooth/strictly-cvx-1.pdf}
{figures/tikz/smooth-nonstrict-convexity.pdf}
{A convex function need not be strictly convex.}
{The highlighted interval is exactly affine, so equality holds in the secant inequality there. The surrounding curve is chosen with nondecreasing slope.}
{smooth-nonstrict-convexity}
\ReviewFigure{The gradient is a vector field}
{figures/optim-smooth/gradient-vector-field.pdf}
{figures/tikz/smooth-gradient-field.pdf}
{The gradient associates a vector in the domain with each evaluation point.}
{The radial field is shown for f(x)=||x||\textsuperscript{2}/2, so the gradient points outward and vanishes at the minimizer. The accompanying bowl is schematic; its axes are not numerically calibrated.}
{smooth-gradient-field}
\ReviewFigure{Stationarity at local extrema}
{figures/optim-smooth/first-order-1.pdf}
{figures/tikz/smooth-stationary-extrema.pdf}
{Vanishing derivative is necessary at differentiable local minima and maxima.}
{The representative double-well function has two minima and a maximum, all with horizontal tangents. Stationarity alone does not distinguish these cases.}
{smooth-stationary-extrema}
\ReviewFigure{A stationary point need not be an extremum}
{figures/optim-smooth/first-order-2.pdf}
{figures/tikz/smooth-stationary-inflection.pdf}
{The first-order necessary condition is not sufficient without convexity.}
{The reconstruction uses f(x)=x\textsuperscript{3} as in the accompanying text and labels the point as a stationary inflection point. The old caption calls the one-dimensional example a saddle point; the more precise terminology is used here.}
{smooth-stationary-inflection}
\ReviewFigure{Stationarity for a convex objective}
{figures/optim-smooth/first-order-3.pdf}
{figures/tikz/smooth-stationary-convex.pdf}
{For differentiable convex functions, a vanishing gradient characterizes global minimizers.}
{The tangent plane touches the bowl at its minimum and supports the graph from below. The plane is horizontal in the depicted three-dimensional coordinates.}
{smooth-stationary-convex}
\ReviewFigure{PCA and least-squares level sets}
{figures/optim-smooth/link-pca.pdf}
{figures/tikz/smooth-pca-quadratic-geometry.pdf}
{Covariance axes and quadratic level-set axes have reciprocal widths.}
{Both panels share the eigenvectors u\textsubscript{1},u\textsubscript{2}. The twelve reflection-symmetric points are centered and have empirical covariance eigenvalues \ensuremath{\lambda}\textsubscript{1}=1 and \ensuremath{\lambda}\textsubscript{2}=0.325\textsuperscript{2}. Their covariance ellipses and the quadratic level sets are derived from these same eigenvalues, with reciprocal axis ratios. The caption inside the drawing specifies the common multiplier used for each family of semiaxes; axis labels follow the covariance normalization C/n in the text.}
{smooth-pca-quadratic-geometry}
\ReviewFigure{A first-order Taylor approximation}
{figures/optim-smooth/taylor-exp-1.pdf}
{figures/tikz/smooth-taylor-line.pdf}
{The affine approximation matches the function and its derivative at the expansion point.}
{The drawn line is the exact tangent to the representative quadratic at x; the vertical gap at z shows the Taylor remainder.}
{smooth-taylor-line}
\ReviewFigure{A tangent plane}
{figures/optim-smooth/taylor-exp-2.pdf}
{figures/tikz/smooth-taylor-plane.pdf}
{The first-order affine approximation of a differentiable function in two dimensions.}
{The tangent plane is computed for the displayed representative paraboloid and touches it at the marked point. The graph and its supporting plane are distinct objects.}
{smooth-taylor-plane}
\ReviewFigure{The gradient is orthogonal to a level curve}
{figures/optim-smooth/level-sets-1.pdf}
{figures/tikz/smooth-level-set-tangent.pdf}
{Differentiating f(c(t))=constant gives a zero pairing with the tangent direction.}
{At the leftmost point of an elliptical level curve, the tangent is vertical and the outward gradient is horizontal. The right-angle marker records orthogonality, not a descent step.}
{smooth-level-set-tangent}
\ReviewFigure{Gradients and nested level sets}
{figures/optim-smooth/level-sets-2.pdf}
{figures/tikz/smooth-nested-level-sets.pdf}
{Gradients point in the direction normal to a level set and toward larger values.}
{The arrows are true normals to elliptical level curves, rather than radial arrows except on the principal axes. All level values increase outward.}
{smooth-nested-level-sets}
\ReviewFigure{Step-size effects in gradient descent}
{figures/optim-smooth/grad-desc-1.pdf}
{figures/tikz/smooth-gradient-step-size.pdf}
{The step size controls progress and oscillation along a quadratic objective.}
{Both trajectories are exact gradient-descent iterates for the displayed quadratic. The larger step oscillates but remains stable because both step sizes lie in (0,2/L), with L=5. The sketch is schematic; no numerical values were legible in it.}
{smooth-gradient-step-size}
\ReviewFigure{Exact line search}
{figures/optim-smooth/grad-desc-2.pdf}
{figures/tikz/smooth-exact-line-search.pdf}
{A differentiable interior line-search optimum makes successive gradients orthogonal.}
{The path uses exact line search for a positive definite quadratic. The right-angle marker is computed from consecutive actual steps; it is not merely a schematic zigzag.}
{smooth-exact-line-search}
\ReviewFigure{The optimal quadratic step size}
{figures/optim-smooth/proof-quadr.pdf}
{figures/tikz/smooth-quadratic-contraction.pdf}
{The operator norm of I-\ensuremath{\tau}C is controlled by the extreme eigenvalues.}
{The highlighted curve is the maximum of the two absolute-value branches. Its minimum occurs at 2/(L+\ensuremath{\mu}), and h(\ensuremath{\tau})<1 exactly for 0<\ensuremath{\tau}<2/L when 0<\ensuremath{\mu}\ensuremath{\leq}L. A representative ratio L/\ensuremath{\mu}=3 determines the drawing; the labels remain symbolic.}
{smooth-quadratic-contraction}
\ReviewFigure{Quadratic upper and lower bounds}
{figures/optim-smooth/up-low-bounds.pdf}
{figures/tikz/smooth-curvature-bounds.pdf}
{Strong convexity and Lipschitz gradient give touching lower and upper quadratic models.}
{Both bounding quadratics have the same value and gradient as f at x. The lower curvature \ensuremath{\mu} and upper curvature L are ordered correctly on both sides of the expansion point.}
{smooth-curvature-bounds}
